% GPT-6 Astra Pro; effort not exposed; /root/replacement; 22 November 2026.
\section[HD-F1-R22: the actual hull-quotient cancellation]{HD-F1-R22: the exact Ext cancellation for a common quotient of rank one}
\label{sec:followup-rm-rank-one-quotient-cancellation}

Retain R21's gluable weakly semiregular adjacent class:
positive virtual rank one, equal-slope semistable torsion-free
endpoints of ranks $(r_A,r_T)=(3,2)$, and any specified
coherent faithful finite point action with an actual lift.
Apply the actual determinant twist in (F1R21.2)--(F1R21.3),
so that $\operatorname{ch}(E)=v_q$ and both endpoint slopes
are zero. Let
\[
 W=\operatorname{Hom}_G(A,T),\qquad h=\dim W\ge2.
 \tag{F1R22.1}
\]
Suppose the actual common-kernel alternative of R21 occurs:
all maps in $W$ factor through its specified rank-one
TF quotient
\[
                p:A\twoheadrightarrow Q,
                       \qquad \mu_H(Q)=0.
 \tag{F1R22.2}
\]
We prove that $h=2$, that $T^{**}$ is equivariantly a
sum of two copies of the actual line $L=Q^{**}$, and
that $Z=T^{**}/T$ must have a nonzero codimension-two
part. More precisely, the actual map
\[
 \operatorname{Ext}^2_G(Z,L)\longrightarrow
                   \operatorname{Ext}^2_G(L^{\oplus2},L)
 \tag{F1R22.3}
\]
has an image of dimension at least three, with the
specified subspace constraint below. In particular
this common-quotient alternative is excluded whenever
$T$ is reflexive, or whenever $T^{**}/T$ has support
of codimension at least three. We do not discard the
codimension-two contribution in the general TF case.

\paragraph{1. Two actual maps identify the target hull.}
At least two of the maps $Q\to T$ induced by $W$ have
independent generic images. Otherwise all their images
would lie in a single saturated rank-one TF subsheaf
$I\subset T$. The nonzero map $Q\to I$ is injective,
so $0=\mu(Q)\le\mu(I)\le0$, using semistability of
$T$ for the last inequality. Double dual gives the
injective map
\[
 \operatorname{Hom}(Q,I)\lhook\joinrel\longrightarrow
                            \operatorname{Hom}(Q^{**},I^{**}).
 \tag{F1R22.4}
\]
Both hulls are line bundles of slope zero. A nonzero
map between them has an effective zero divisor of
degree zero, hence is an isomorphism; its Hom space
has dimension one. This would give $h\le1$, a
contradiction. The injectivity in (F1R22.4) follows
because the target $I$ is torsion free, so a map
vanishing generically is zero.

Choose two invariant maps with independent generic
images, and double-dualize them. They give the actual
map
\[
                      d:L^{\oplus2}\longrightarrow T^{**},
                                 \qquad L=Q^{**}.
 \tag{F1R22.5}
\]
It is generically an isomorphism. The determinant
section has degree zero, since $\mu(T^{**})=\mu(L)=0$,
so it has no zero divisor. On the common big open
set where both sheaves are locally free, it is
therefore an isomorphism. The two sheaves are reflexive;
the actual restriction/pushforward inverse across
codimension at least two extends the isomorphism
and its inverse to all of $X$. Thus $d$ is an
isomorphism globally.

The maps chosen in (F1R22.5) are invariant, and the
quotient and double-dual operations carry their
induced coherent lifts. Consequently $d$ is an
isomorphism for the specified action, not merely
an ordinary bundle isomorphism. Identify $T^{**}$
with $L^{\oplus2}$ by this very map, and put
\[
       q:A\xrightarrow{p}Q\longrightarrow L,\qquad
       j:T\lhook\joinrel\longrightarrow L^{\oplus2},\qquad
       Z=\operatorname{coker}j.
 \tag{F1R22.6}
\]
Every invariant forward map $f:A\to T$ now satisfies
\[
                           j f=v_f q,
                                   \qquad v_f\in\mathbf C^2.
 \tag{F1R22.7}
\]
Indeed it factors through $Q$, and its double dual
is a map $L\to L^{\oplus2}$, given by a constant
column. Equality on the generic point proves
(F1R22.7) globally because the target is torsion free.
In particular $W$ injects into $\mathbf C^2$, so
$h=2$. Write $j_1,j_2:T\to L$ for the two actual
components of $j$.

\paragraph{2. The balanced kernel and its pure-line image.}
Let $K_E=\ker c^2_{v_q}$ be the normalized balanced
kernel. Weak semiregularity, via the exact identity-sector
comparison of R05, gives
$\operatorname{ev}_E^2(K_E)=0$. The second actual
endpoint Hom sequence gives, for every $k\in K_E$,
a map $f_k\in W$ such that
\[
                       \operatorname{ev}_T^2(k)=f_k[2]\delta,
                        \qquad \delta:T\longrightarrow A[2].
 \tag{F1R22.8}
\]
Only this precise Yoneda kernel membership is used;
no choice of $f_k$ linear in $k$ is required.

For the actual line bundle $L$ put
\[
 \lambda(k)=\operatorname{ev}_L^2(k)\in H^2(X,\mathcal O_X),
 \qquad \Lambda=\lambda(K_E),\qquad r=\dim\Lambda.
 \tag{F1R22.9}
\]
The pure-line full trace is multiplication by
$e^{c_1(L)}$, hence is injective on $H^2(\mathcal O_X)$.
The complete rank-four quotient calculation in R10
therefore gives
\[
                                  2\le r\le6.
 \tag{F1R22.10}
\]
Explicitly, with
$\pi_\ell(v,\eta)=\eta+\iota_v\ell$, $\ell=c_1(L)$,
the character map is
$e^\ell\bigwedge^2\pi_\ell$, and
$\operatorname{rank}(\bigwedge^2\pi_\ell|_{K_E})\ge2$.
This is the proved pure-character quotient bound,
with the exact divisor shear retained.

Naturality and the actual direct-sum isomorphism
(F1R22.5) give
\[
 \operatorname{ev}_{L^{\oplus2}}^2(k)
             =\lambda(k)1_{\mathbf C^2},\qquad
 \operatorname{ev}_{L^{\oplus2}}^2(k)j
             =j[2]\operatorname{ev}_T^2(k).
 \tag{F1R22.11}
\]
Set the fixed actual class
\[
                          \beta=q[2]\delta
                              \in\operatorname{Ext}^2_G(T,L).
 \tag{F1R22.12}
\]
Combining (F1R22.7)--(F1R22.12) yields
\[
       \bigl(\lambda(k)j_1,\lambda(k)j_2\bigr)
                 =\bigl(v_{f_k,1}\beta,v_{f_k,2}\beta\bigr).
 \tag{F1R22.13}
\]
Thus every row class $\lambda(k)j_a$ lies in the
single fixed line $\mathbf C\beta$, including the
possibility $\beta=0$.

The induced action on $\operatorname{Ext}^2(L,L)$
identifies with translation on $H^2(\mathcal O_X)$,
which is the identity; the line tensor cancels in
endomorphisms. The two selected invariant columns
in (F1R22.5) also make the multiplicity space trivial
for this action. Hence
\[
 \Lambda\otimes(\mathbf C^2)^*
          \subset\operatorname{Ext}^2_G(L^{\oplus2},L)
 \tag{F1R22.14}
\]
is a subspace of dimension $2r$ for the actual lift.

\paragraph{3. The exact codimension-two cancellation map.}
Apply equivariant Hom into $L$ to the actual exact
sequence $0\to T\xrightarrow{j}L^{\oplus2}\to Z\to0$.
The relevant terms are
\[
 \operatorname{Ext}^2_G(Z,L)
       \xrightarrow{\pi^*}\operatorname{Ext}^2_G(L^{\oplus2},L)
       \xrightarrow{j^*}\operatorname{Ext}^2_G(T,L).
 \tag{F1R22.15}
\]
By (F1R22.13), the restriction of $j^*$ to the
$2r$-dimensional subspace (F1R22.14) has image of
dimension at most one. Exactness therefore gives
\[
 \dim\left(\operatorname{im}\pi^*
        \cap\bigl(\Lambda\otimes(\mathbf C^2)^*\bigr)\right)
                                    \ge2r-1\ge3.
 \tag{F1R22.16}
\]
When $\beta=0$, the stronger bound is $2r$.
This proves the asserted lower bound for the actual
map (F1R22.3), not just for the dimension of a proposed
source vector space. In particular
\[
 \operatorname{ext}^2_G(Z,L)\ge3,
       \qquad \dim H^2(X,Z\otimes L^{-1})^G\ge3,
 \tag{F1R22.17}
\]
where the second bound uses invariant
CY4 Serre duality
$\operatorname{Ext}^2_G(Z,L)^*\simeq
\operatorname{Ext}^2_G(L,Z)$, followed by the actual
line-bundle identity
$\operatorname{Ext}^2(L,Z)=H^2(Z\otimes L^{-1})$.
The superscript on the last cohomology means its
induced action through these actual equivariant maps.

For clarity, a sheaf supported in codimension at least
three on the smooth fourfold has
$\mathcal Ext^i(Z,L)=0$ for $i<3$. Locally this follows
from a regular sequence of length three in its
annihilator in the regular local ring: the successive
change-of-rings Ext sequences shift the first possible
Ext degree by three. The local-to-global Ext spectral
sequence then gives $\operatorname{Ext}^2(Z,L)=0$.
This is also the precise vanishing in Huybrechts--Lehn,
\emph{The Geometry of Moduli Spaces of Sheaves}, second
edition, Proposition~1.1.6(i), applied after the actual
line twist by $L^{-1}$; its complete published proof
was checked in the official publisher excerpt.
Thus (F1R22.17) forces a nonzero codimension-two part
of $Z$; $Z$ already has codimension at least two
because it is the quotient of the reflexive hull by the TF sheaf.
This proves the claimed exclusions, including every
reflexive target in the common-quotient alternative.

\paragraph{The corrected index and exact remaining scope.}
A provisional attempted strengthening used only
$\operatorname{Ext}^1(Z,L)=0$ at codimension two to
claim injectivity of $j^*$ in degree two. The full
sequence (F1R22.15) shows precisely why that inference
fails: the preceding term is $\operatorname{Ext}^2(Z,L)$.
The attempted general TF exclusion was rejected before
promotion and is recorded separately in the checking
log. Its correct replacement is the actual cancellation
bound (F1R22.16), which retains that term.

The underlying characteristic/trace and pure-line
quotient maps are the previously checked published
Buchweitz--Flenner framework and the explicit R10
linear algebra. The double-dual maps use the same
actual big-open extension as R18. The proof does
not compute or assert the nonexistence of the
codimension-two cancellation when it is allowed.
It narrows the common-quotient alternative to
$h=2$ and the indicated singular target with a
specified nonzero Ext image. The common-image
alternative of R21 and the remaining cancellation
problem are not replaced by an assumed bundle or
by a hypothetical numerical character.
